\documentclass[12pt]{article}
\usepackage{color,soul}
% Import preambles and macros for notes
\input{../../../../preambles/notes-preamble.tex}
\input{../../../../preambles/notes-macros.tex}
\input{../../../../preambles/theorem-system-section.tex}

\title{MATH 420 Lecture 18}
\author{Deepak Jassal}
\date{March 18\textsuperscript{th}, 2026}

\begin{document}
\maketitle
\stepcounter{section}
\propp{\textit{Prime Avoidance}\\
    Let $\mathcal{P}_1,\dots,\mathcal{P}_r$ be prime ideals in a ring $R$. If $I\subsetneq \mathcal{P}_i$ ($I$ an ideal of $R$) for $i=1,\dots,r$, then $I\subsetneqq\bigcup_{j=1}^r\mathcal{P}_i$.
}{
    If $r=1$, the claim is obvious. Hence, we assume that $r\geq2.$ Suppose that $I\subset I\subsetneqq\bigcup_{j=1}^r\mathcal{P}_i$, but $I$ is not contained in the union of fewer $\mathcal{P}_i$'s. In particular for $1\leq i\leq r$ $I\subsetneqq\bigcup_{j\neq i}\mathcal{P}_j$. Hence there exists $a_i\in I\setminus\bigcup_{j\neq i}\mathcal{P}_j$. But then $a_i\in\mathcal{P}_i$. Take $a=a_1\cdots a_{r-1}+a_r\in I$.\\
    \textit{Claim:} $a\not\in\mathcal{P}_i$ for $1\leq i\leq r$. To see this note that $a_1,\dots,a_{r-1}\not\in\mathcal{P}_r$.\\
    Since $\mathcal{P}_r$ is prime, their product must not lie in $\mathcal{P}_r$. By definition, $a_r\in\mathcal{P}_r$. Hence $a\not\in\mathcal{P}_r$. If $1\leq i\leq r-1$, then $a_1\cdots a_{r-1}\in\mathcal{P}_i$, but $a_r\not\in\mathcal{P}_i$. $a\in I$, but $a\not\in\mathcal{P}_i$ for $1\leq i\leq r$, contradicts $I\subset\bigcup_{j=1}^r\mathcal{P}_j$.
}
$\Z[\sqrt{-5}]$ is a quadratic extention of the $\Z$ (the ring of integers extended with roots for unsolvable (monoic) quadratic polynomials). $(5)\in\Z$ is irreducible. But $(5)=(\sqrt{-5})\in\Z[\sqrt{-5}]$ is reducible.\\
Let $\phi:A\to B$ be a ring homomorphism ($B$ has the structure of an $A$-algebra). If $b$ is an ideal in $B$, then $\phi^{-1}b$ is an ideal in $A$. 
\lemp{}{
    If $b$ is an ideal in $B$, then $\phi^{-1}b$ is an ideal in $A$.\\
}{
    \textit{Recall:}
    \[
        \phi^{-1}b=\{a\in A:\phi(a)\in b\},
    \]
    $0_A\in\phi^{-1}(\{0_B\})$, so we need to check 
    \begin{enumerate}
        \item $0_B\in b$
        \item If $a_1,a_2\in\phi^{-1}(b)$ then there exists $b_1,b_2\in b$ such that $a_i\in\phi^{-1}(\{b_i\})$, (i.e., $b_i=\phi(\{a_i\})\in b$). $b_1+b_2\in b$, $b$ is an ideal. $b_1+b_2\phi(a_1)+\phi(a_2)=\phi(a_1+a_2)\Rightarrow a_1+a_2\in\phi^{-1}(b)$. Supose $r\in A$, $a\in\phi^{-1}(b)$. Then $\phi(ra)=\phi(r)\phi(a)$. $\phi(a)\in b$, hence $\phi(r)\phi(a)\in b$. Hence $\phi(ra)\in b$ and $ra\in\phi^{-1}(b)$.
    \end{enumerate}
    Thus, we have our result.
}
\defn{Ideal Contractions}{
    Let $\phi:A\to B$, be a ring homomorphism. For an ideal $b$ of $B$, $\phi^{-1}(b)$ is called the \underline{contraction} of $b$ in $A$.
}
\ex{
    $\Z[\sqrt{-5}]$ contracts to what in $\Z$?\\
    $\phi^{-1}(\Z[\sqrt{-5}])=(6)$.
}
In general $\phi(I)$ is \underline{not} an ideal in $B$. But we can take the ideal generated by $\phi(I)$ and call that the \underline{extension} of $I$.
\lemp{}{
    If $A$ is a subring of $B$, and $b\subset B$ an ideal ($b^c$ is the contraction of $b$ in $A$), then $b^c=b\cap A$ ($\phi:A\to A$ is the identity map).
}{
    $a\in b^c=\phi^{-1}(b)$ \textit{iff} $\phi(a)\in b$, but $\phi(a)=a$, so $a\in b$. Hence $b^c=b\cap A$.
}
\section*{Chinese Remainder Theorem}
\textit{Classical:} If $m_1,\dots,m_k$ are pairwise coprime integers. and $a_1,\dots,a_k$ are arbitrary integers, then there exists $x\in\Z_{\geq1}$ such that $x-a_i$ is divisible by $m_i$ for $i=1,\dots,k$.\\
We have to interpret the second half of the statement as a statement about the map $\phi:R\to R/m_1\times\cdots R/m_k$, $a\mapsto(a+m_1,\dots,a+m_k)$.
\end{document}